\documentclass[11pt]{article}
\usepackage[margin=1in]{geometry}
\usepackage{amsmath,amssymb,amsthm,mathtools}
\usepackage[colorlinks=true,linkcolor=blue,citecolor=blue,urlcolor=blue]{hyperref}
\newtheorem{theorem}{Theorem}[section]
\newtheorem{lemma}[theorem]{Lemma}
\newtheorem{proposition}[theorem]{Proposition}
\newtheorem{corollary}[theorem]{Corollary}
\theoremstyle{remark}
\newtheorem{remark}[theorem]{Remark}
\newcommand{\Z}{\mathbb Z}
\newcommand{\Q}{\mathbb Q}
\newcommand{\R}{\mathbb R}
\newcommand{\N}{\mathbb N}
\newcommand{\rk}{\operatorname{rank}_{\Z}}
\newcommand{\ark}{\widetilde{\operatorname{rank}}_{\Z}}
\newcommand{\res}{\mathrel{\geq_{\Z}}}
\newcommand{\degto}{\mathrel{\rightsquigarrow_{\Z}}}
\DeclareMathOperator{\Norm}{Norm}
\title{Shared-leg rectangular matrix multiplication\protect\\over arbitrary rings}
\author{Yang Zhou\textsuperscript{1,2,3}\qquad Mingsheng Wang\textsuperscript{2}\\[0.8ex]
  {\small\textsuperscript{1}State Key Laboratory of Information Security}\\
  {\small\textsuperscript{2}Institute of Information Engineering, Chinese Academy of Sciences}\\
  {\small\textsuperscript{3}School of Cyber Security, University of Chinese Academy of Sciences}\\\\[0.8ex]
  {\small\href{mailto:zhouyang2000@iie.ac.cn}{\texttt{zhouyang2000@iie.ac.cn}}}\\
  {\small\href{mailto:wangmingsheng@iie.ac.cn}{\texttt{wangmingsheng@iie.ac.cn}}}}
\date{}
\begin{document}
\maketitle
\begin{abstract}
We give an integral extension of the rectangular bound in
\href{https://arxiv.org/abs/2610.10494}{arXiv:2610.10494v1}.
The new algebraic ingredients are an integer asymptotic-rank bound for
convolution and an integer replacement for Fourier separation.
Together with a discrete proof of the real-variable profile implication,
they produce integer bilinear algorithms valid over every associative
ring, including noncommutative rings and rings without identity.
All bounds concern arithmetic operations, with arbitrary positive
slack in the exponent; no bit-complexity bound is asserted.
\end{abstract}

\section{Statement and computational model}

For $k\geq0$ set
\[
 f(k)=\begin{cases}
 2,&0\leq k\leq\tfrac12,\\
 1+k+\dfrac1{4k},&k>\tfrac12.
 \end{cases}
\]
\begin{theorem}[Universal ring version]\label{thm:ring}
For every $k\geq0$ and $\varepsilon>0$ there are a family of integer
bilinear circuits and a constant $C=C(k,\varepsilon)>0$ such that,
for every associative ring $R$ and every integer $n\geq1$, the circuit
for size $n$ multiplies an $n\times\lceil n^k\rceil$ matrix by a
$\lceil n^k\rceil\times n$ matrix using at most
$Cn^{f(k)+\varepsilon}$ additions, negations, and multiplications.
The circuits and $C$ are chosen before $R$. Commutativity, an identity,
and restrictions on characteristic or zero divisors are unnecessary.
\end{theorem}

Here an integer bilinear algorithm forms integer linear combinations
of the first input, integer linear combinations of the second input,
multiplies each first combination on the left of its second
combination, and recombines the products with integer coefficients.
Integer multiples are implemented by additions and negations.
For $c\in\Z$ and $r\in R$, $cr$ means the integer action on the
additive group: repeated addition, with negation for negative $c$.
It requires neither an element $c$ in $R$ nor a multiplicative identity.
Each addition, negation, or ordered multiplication costs one operation.
Input reads, zero, copying and reindexing available values, and choosing
output coordinates cost no arithmetic operations. Finite intermediate
families are computed once and may be reused. There is no primitive
integer scalar multiplication: it is compiled into additions and negations.
The positive exponent slack absorbs logarithmic overhead; the theorem
does not assert $O(n^{f(k)})$ without slack.

\paragraph{Background used.}
For background on bilinear algorithms and block recursion, consult
\cite[Chapters 14--15]{BCS}; for finite fields, \cite[Chapter 2]{LN}.
We use elementary finite free module and determinant algebra,
the existence of finite fields of every prime-power order, and standard
real analysis (including Fekete's lemma and Stirling's estimate).
Its only specialized external theorem is abstract spectral rank duality,
stated below with its hypotheses. The integral constructions and discrete
parameter argument are proved here; spectral duality itself is cited
rather than reproved. A declaration-level correspondence with the checked
Lean development is in \path{formalization/proof-correspondence.md}.

\paragraph{Proof route.}
An integer coefficient identity gives an algorithm over every ring
because its two multiplication factors retain their order. It therefore
suffices to obtain one fixed integer base algorithm with the required
rank-to-size ratio. Integral spectral duality reduces that task to
bounding character values of matrix tensors. Coefficient extraction and
convolution supply the shared-leg entropy inequalities; the discrete
profile argument then bounds the character parameters by a parabola.
The final step selects the base algorithm and charges all coefficient
arithmetic in its recursion.

\section{An integral spectrum with the right order}

An integer tensor is a finite three-index array $T_{ijh}$, written
as the formal expression
\[
 T(x,y,z)=\sum_{i,j,h}T_{ijh}x_i y_j z_h.
\]
The three variable groups are its \emph{legs}. These variables record
coefficients; they do not impose commutativity on the ring where the
algorithm will run. A rank-one tensor is $L(x)M(y)N(z)$, with integer
linear forms. Thus $T=\sum_{\ell=1}^rL_\ell(x)M_\ell(y)N_\ell(z)$
describes $r$ ordered products: evaluate $L_\ell$ on the first input
and $M_\ell$ on the second, multiply in that order, and use the
coefficients of $N_\ell$ to recombine the outputs. Tensor rank counts
these products; the final section also charges additions and negations.

A restriction independently substitutes integer linear forms on each
leg. Direct sum uses disjoint coordinate groups on all three legs,
with every mixed coefficient zero. Tensor product pairs coordinate
indices and multiplies the corresponding integer coefficients.
All coordinate spaces below are finite free $\Z$-modules.
Write $A\res B$ when $A$ restricts to $B$. For $m\in\N$, let
$D_m=\sum_{i=1}^m x_i y_i z_i$ and $1=D_1$.
The integer rank $\rk(T)$ is the least $m$ with $D_m\res T$.

We work in the semiring $\mathcal S$ consisting of $0$ and the integer
tensors that restrict to $1$, modulo mutual integer restriction.
Addition is full direct sum and multiplication is tensor product.
This choice matters: all nonzero integer tensors would be unsuitable,
since $2xyz$ cannot restrict to $xyz$ over $\Z$.

The semiring is closed under both operations and under leg
permutations. Every nonzero $T\in\mathcal S$ satisfies
\[
 D_{\rk(T)}\res T\res1.
\]
Expansion in a basis proves finiteness of rank. The diagonal tensors
embed $\N$ with its usual order: the reverse implication follows by
flattening after base change to $\Q$. Restrictions respect both
operations. Finally, for any $A$ and nonzero $B$, some $m$ satisfies
$mB\res A$, because $B\res1$ and $D_{\rk(A)}\res A$.
Thus restriction is a Strassen preorder on $\mathcal S$.

Here the preorder convention for the abstract theorem is
$B\preceq A$ if $A\res B$. The theorem applies to a commutative
semiring with a preorder that induces the usual order on $\N$,
respects addition and multiplication, and has the cofinality property
$A\preceq mB$ for some $m\in\N$ whenever $B\ne0$.
It asserts that the normalized monotone semiring homomorphisms form a
nonempty spectrum and, when $1\preceq T$, their supremum is the
asymptotic rank of $T$. The preceding paragraph verifies each hypothesis;
commutativity here concerns tensor classes, independently of the eventual
coefficient ring of the algorithm.

Let $\Delta(\mathcal S)$ be the set of normalized, additive,
multiplicative, restriction-monotone maps $\lambda:\mathcal S\to\R_{\geq0}$.
Set $\ark(0)=0$ and, for nonzero $T\in\mathcal S$, define
\[
 \ark(T)=\exp\left(\inf_{m\geq1}
                    \frac{\log\rk(T^{\otimes m})}{m}\right).
\]
Rank is submultiplicative under tensor product, so the logarithms are
subadditive. Fekete's lemma identifies this infimum with the limit in
the next display. This infimum definition is the one used in Lean.
The abstract spectral rank theorem~\cite[Theorem 2.13,
Corollary 2.14, Theorem 2.16]{Zui}
therefore gives, for nonzero $T\in\mathcal S$,
\begin{equation}\label{eq:spectrum}
 \ark(T)=\lim_{m\to\infty}\rk(T^{\otimes m})^{1/m}
 =\sup_{\lambda\in\Delta(\mathcal S)}\lambda(T).
\end{equation}
The nonzero case has $T\res1$, so every character has value at least
$1$, as required by that theorem. For $T=0$, both rank and every
character value are zero, and the same display holds directly.
We call these maps characters. In particular,
$1\leq\lambda(T)\leq\rk(T)$ for nonzero objects, and
$\lambda(T)\leq\ark(T)$ by applying the rank bound to tensor powers.
Thus a character is a numerical size estimate compatible with independent
tasks, tensor products, and integer preprocessing. Each character gives
a rank lower bound; spectral duality says that bounding all of them also
bounds asymptotic rank.

For a positive integer $q$, the oriented dot products are
\[
 B_X(q)=x\sum_{j=1}^q y_jz_j,\quad
 B_Y(q)=y\sum_{j=1}^q x_jz_j,\quad
 B_Z(q)=z\sum_{j=1}^q x_jy_j.
\]
They belong to $\mathcal S$. For each character, multiplicativity and
monotonicity give parameters $p_i=\log_2\lambda(B_i(2))$, $i\in\{X,Y,Z\}$,
with
\begin{equation}\label{eq:dot}
 \lambda(B_i(q))=q^{p_i},\qquad 0\leq p_i\leq1.
\end{equation}
Comparison of $q^m$ with adjacent powers of $2$ proves the power law.
For positive integers $a,b,c$, define the matrix multiplication tensor
\[
 \langle a,b,c\rangle=\sum_{i=1}^a\sum_{j=1}^b\sum_{h=1}^c
             x_{ij}y_{jh}z_{hi}.
\]
Its factorization into the three oriented dot products gives
\begin{equation}\label{eq:mm}
 \lambda(\langle a,b,c\rangle)=a^{p_Y}b^{p_Z}c^{p_X}.
\end{equation}

\section{Coefficient extraction replaces division}

Write $A\degto S$ when integer Laurent-polynomial leg maps transform
$A$ to a tensor family $F(t)=t^\ell S+\text{terms of degree greater than }\ell$
for some $\ell\in\Z$. This also allows $S=0$.
In the weighted constructions below, a variable of integer weight $w$
is replaced by $t^w$ times that variable. The lowest total-weight terms
give the leading coefficient. The parameter $t$ is a formal bookkeeping
device; coefficient extraction makes the resulting comparison integral.

\begin{lemma}[Integral degeneration monotonicity]\label{lem:deg}
If $A,S\in\mathcal S$ and $A\degto S$, then
$\lambda(A)\geq\lambda(S)$ for every character.
\end{lemma}
\begin{proof}
Shift the leg maps so their entries are polynomials and the shifted
$\ell$ is nonnegative. Write the resulting tensor as
$F(t)=t^\ell S+\text{higher-degree terms}$, $\ell\geq0$.
For an integer $d\geq0$, the coefficient tensor on three copies of
$\Z^{d+1}$ is
\[
 H_d=\sum_{\substack{i,j,h\geq0\\i+j+h=d}}
            u_i v_j w_h.
\]
It has a coefficient $1$ and at most $(d+1)^2$ terms, so it lies in
$\mathcal S$ and has rank at most $(d+1)^2$.
For each $m\geq1$, extracting degree $m\ell$ from the powered leg maps
gives the restriction
\[
 A^{\otimes m}\otimes H_{m\ell}\res S^{\otimes m},
\]
explicitly: if $L_X(t)$ is the powered first-leg map, send
$x\otimes u_i$ to $[t^i]L_X(t)x$, and use the analogous maps on the
other legs. The condition $i+j+h=m\ell$ extracts exactly that
coefficient of $F(t)^{\otimes m}$, which is $S^{\otimes m}$.
Thus
\[
 \lambda(S)^m\leq(m\ell+1)^2\lambda(A)^m.
\]
The extra factor tends to $1$ after taking $m$th roots. Here $\ell$
is fixed by the original degeneration, independently of $m$.
\end{proof}

\section{The integral convolution bound}

For positive $r,s$ define
\[
 C(r,s)=\sum_{i=0}^{r-1}\sum_{j=0}^{s-1}x_i y_j z_{i+j},
 \qquad N=r+s-1.
\]
Its coefficient of $x_0y_0z_0$ is $1$, so $C(r,s)\in\mathcal S$.
The later separator needs only $\lambda(C(r,s))\leq N$. A single
scaled interpolation identity cannot give this over the integers:
its multiplier must also be eliminated. We construct finitely many
identities whose multipliers can be patched by B\'ezout.

\begin{lemma}\label{lem:conv}
For every fixed $r,s\geq1$, with $N=r+s-1$, there is a constant
$K_{r,s}>0$ such that
\[
 \rk(C(r,s)^{\otimes m})\leq K_{r,s}N^m\quad(m\geq1).
\]
Consequently every character satisfies $\lambda(C(r,s))\leq N$.
\end{lemma}
\begin{proof}
Let $B=\Z[t]/(g)$, where $g$ is monic of degree $d\geq1$; $B=\Z$
with $d=1$ is also allowed. Its standard basis starts with $1$.
For any nodes $a_0,\ldots,a_{N-1}\in B$, form the Vandermonde matrix
$V_{\ell k}=a_\ell^k$ ($0\leq\ell,k<N$), and put
$\delta=\det V$. The identity $\operatorname{adj}(V)V=\delta I$ gives
\[
 \delta C(r,s)=\sum_{\ell=0}^{N-1}
 \left(\sum_{i=0}^{r-1}a_\ell^i x_i\right)
 \left(\sum_{j=0}^{s-1}a_\ell^j y_j\right)
 \left(\sum_{k=0}^{N-1}\operatorname{adj}(V)_{k\ell}z_k\right).
\]
Define $E=\Norm_B(\delta)\in\Z$ as the determinant of multiplication
by $\delta$ on the free integer module $B$. Its adjugate supplies
$\gamma\in B$ with $\gamma\delta=E\cdot1_B$: if $M_\delta$ is that
multiplication matrix, take the coordinates of $\gamma$ to be
$\operatorname{adj}(M_\delta)$ times the coordinates of $1_B$.
The identity $M_\delta\operatorname{adj}(M_\delta)=EI$ proves the claim.
Multiplying the output
forms above by $\gamma$ therefore gives a length-$N$ decomposition of
$EC(r,s)$ over $B$. Both adjugate identities hold without inversion.

Take the $m$th tensor power \emph{in this same algebra}. It decomposes
$E^mC(r,s)^{\otimes m}$ with $N^m$ terms. To descend it once to integers,
let $\pi:B\to\Z$ extract the coefficient of $1$. For $u,v,w\in B$,
write $u=\sum_{a<d}u_at^a$ and $v=\sum_{b<d}v_bt^b$. Then

\[
 \pi(uvw)=\sum_{0\leq a,b<d}u_a v_b\,\pi(t^{a+b}w).
\]
Since $\pi(1)=1$, this produces an integer decomposition with at most
$d^2N^m$ terms. Only $\Z$-linearity of $\pi$ is used. Descending after
powering keeps the overhead $d^2$ independent of $m$.

First choose $B_0=\Z$ and the nodes $0,1,\ldots,N-1$.
The resulting Vandermonde determinant is a nonzero integer $E_0$.
For each prime $p\mid E_0$, choose $d_p\geq1$ with $p^{d_p}\geq N$
and a monic integer lift $g_p$ of an irreducible polynomial of
degree $d_p$ over $\mathbb F_p$.
The algebra $B_p=\Z[t]/(g_p)$ reduces to $\mathbb F_{p^{d_p}}$.
Choose $N$ nodes with distinct residues. The Vandermonde determinant
is nonzero in that residue field; its norm, reduced modulo $p$, is
therefore nonzero. Thus the integer multiplier $E_p$ satisfies
$p\nmid E_p$.

Let $I=\{0\}\sqcup\{p:p\text{ is prime and }p\mid E_0\}$ and $d_0=1$.
The multipliers $(E_i)_{i\in I}$ have gcd $1$, as do their $m$th powers.
For each $m\geq1$, choose integers $b_{i,m}$ such that
\[
 \sum_{i\in I}b_{i,m}E_i^m=1.
\]
Multiply the descended rank decompositions by $b_{i,m}$ in their
third legs and add them. This gives an exact integer decomposition
of $C(r,s)^{\otimes m}$ whose number of terms is at most
\[
 \left(\sum_{i\in I}d_i^2\right)N^m.
\]
All algebras, nodes, and dimensions were fixed before $m$; hence
$K_{r,s}=\sum_{i\in I}d_i^2$ is a fixed positive constant. Applying a
character and taking $m$th roots gives the required bound $\lambda(C(r,s))\leq N$.
\end{proof}

\begin{remark}
The output flattening over $\Q$ has rank $N$, and its tensor powers
have rank $N^m$, so the bound also gives $\ark(C(r,s))=N$.
The auxiliary algebras appear only in the proof of an integer identity.
For example, $C(2,2)$ uses the ordinary nodes $0,1,2$ with multiplier
$E_0=2$, and the nodes $0,1,t$ in $\Z[t]/(t^2+t+1)$ with multiplier
$E_2=3$. Thus the powers of $2$ and $3$ can be combined by B\'ezout.
The second multiplier removes the prime obstruction left by the first.
\end{remark}

\section{An integral shared-leg separator}

Character additivity applies to full direct sums. Shared-$X$ sectors
must therefore acquire independent $X$ labels. The convolution below
tracks the difference of two sector labels, and a quadratic weight
retains precisely the terms where they agree.

\begin{proposition}\label{prop:sep}
Let $M\geq1$, and let $X,Y_h,Z_h$ be finite free coordinate modules.
Suppose $A=\sum_{h=1}^M A_h$, with
$A_h\in X\otimes Y_h\otimes Z_h$, where the $Y_h$ are disjoint,
the $Z_h$ are disjoint, and $A_h\in\mathcal S$.
There is an integer restriction followed by an integral degeneration
\[
 A\otimes C(2M-1,2M-1)\degto
 \bigoplus_{h=1}^M\bigl(A_h\otimes B_X(M)\bigr).
\]
The $X$-spaces on the right are independent copies.
\end{proposition}
\begin{proof}
Choose coordinate bases indexed by $a$ in $X$, and by $b,c$ in
$Y_h,Z_h$. Let $c_{habc}$ be the coefficient of $x_a y_{hb}z_{hc}$ in $A$.
The convolution coordinates $\xi_i,\eta_j,\zeta_k$ have ranges
$0\leq i,j\leq2M-2$ and $0\leq k\leq4M-4$.
Introduce target variables $X_{ag},Y_{hbu},Z_{hcv}$,
where $g,u,v\in\{1,\ldots,M\}$.
Use the following three substitutions, with every unmatched
coordinate sent to zero:
\begin{align*}
 x_a\xi_i&\longmapsto\sum_{g:\ i=2g-2}X_{ag},\\
 y_{hb}\eta_j&\longmapsto\sum_{u:\ j=u-h+M-1}Y_{hbu},\\
 z_{hc}\zeta_k&\longmapsto\sum_{v:\ k=v+h+M-3}Z_{hcv}.
\end{align*}
All indicated indices lie in the convolution ranges.
The support equation $i+j=k$ becomes
\[
 u-v+2(g-h)=0.
\]
Give the target variables the integer weights
\[
 w(X_{ag})=g^2,\qquad
 w(Y_{hbu})=hu-h^2,\qquad
 w(Z_{hcv})=-hv.
\]
On the supported terms their sum is $(g-h)^2$.
The lowest-weight terms consequently have $g=h$ and $u=v$,
and are exactly
\[
 \sum_{h,a,b,c,u}c_{habc}X_{ah}Y_{hbu}Z_{hcu}.
\]
This is the stated full direct sum. Every substitution has integer
coefficients; no roots of unity or normalization denominators occur.
To obtain the nonnegative polynomial exponents used in Lean, replace
the three weights by $w_X,w_Y+M^2,w_Z+M^2$. Their supported sum is
$2M^2+(g-h)^2$, so the same target is the coefficient of degree $2M^2$.
\end{proof}

\begin{corollary}[Integral shared-leg entropy]\label{cor:entropy}
Let $s\geq1$ and $T=\sum_{i=1}^sT_i$ have the matched-sector structure
of Proposition~\ref{prop:sep}, with $T_i\in\mathcal S\setminus\{0\}$.
For every character and every probability vector
$q=(q_1,\ldots,q_s)$, $q_i\geq0$, $\sum_iq_i=1$,
\[
 \lambda(T)\geq
 \exp(p_X H(q))\prod_i\lambda(T_i)^{q_i},
 \qquad H(q)=-\sum_i q_i\log q_i,
\]
where $0\log0=0$. For $A\in\mathcal S$ with $A\degto T$, the same
lower bound holds with $\lambda(A)$ on the left.
\end{corollary}
\begin{proof}
Deleting every sector except one restricts $T$ to a nonzero $T_i$,
so $T$ itself admits the unit. In a tensor power, the second and third
legs retain their sector-word labels, whereas the first leg is shared.
For rational $q$, let $m\to\infty$ through multiples of a common
denominator. Retain in $T^{\otimes m}$ only matched words with
$mq_i$ occurrences of sector $i$, by deleting all other second- and
third-leg coordinates. Different retained words have zero mixed
coefficients, so they are matched sectors. There are
$M=m!/\prod_i(mq_i)!$ sectors, each having character value
$\prod_i\lambda(T_i)^{mq_i}$.
Lemmas~\ref{lem:deg}, \ref{lem:conv} and Proposition~\ref{prop:sep}
give
\[
 (4M-3)\lambda(T)^m
 \geq M^{1+p_X}\prod_i\lambda(T_i)^{mq_i}.
\]
Taking $m$th roots leaves only
\[
 \lambda(T)\geq M^{p_X/m}\prod_i\lambda(T_i)^{q_i}
                    \left(\frac{4M-3}{M}\right)^{-1/m}.
\]
Since
$1\leq(4M-3)/M\leq4$ and
$\log M=mH(q)+O_s(\log(m+1))$, the result follows.
The last estimate is Stirling's formula applied to the factorials,
with $0!=1$ for zero counts.
Continuity extends it to all probability vectors.
For each fixed $m$ the degeneration has already been handled by
Lemma~\ref{lem:deg}; no uniform degree estimate in the type-counting
limit is needed.
\end{proof}

\section{Integer sector constructions}

The entropy inequality now applies to two integer constructions: a
filtration supplies concavity, and a three-sector degeneration couples
the two profiles. The orientations used below are $C_{XY}=C$ and
$C_{XZ}(r,s)=\sum_{i<r,j<s}x_i z_j y_{i+j}$, with nonnegative indices.

\paragraph{Determinant filtration.}
Fix positive integers $r,s$. Let $\mathcal V_e=\Z[u,v]_e$ be the free
module of homogeneous binary forms of degree $e\geq0$, and let
$\sigma,\tau$ be a basis of an auxiliary rank-two module. In
$\mathcal V_e\otimes\langle\sigma,\tau\rangle$, use
\[
 A_i=u^{e-i}v^i\sigma,\quad B_i=u^{e-i}v^i\tau
       \quad(0\leq i\leq e),\qquad
 K_j=B_j-A_{j+1}\quad(0\leq j<e).
\]
The basis $(A_0,\ldots,A_e,B_e,K_0,\ldots,K_{e-1})$ is unimodular:
its inverse is $B_j=K_j+A_{j+1}$. Evaluation
$\sigma\mapsto u$, $\tau\mapsto v$ sends the first $e+2$ vectors to
the monomial basis of $\mathcal V_{e+1}$. For $e\geq1$, the kernel is
$(u\tau-v\sigma)\mathcal V_{e-1}$; for $e=0$, it is zero.

In monomial bases, multiplication of forms with $r$ and $s$ coefficients
is exactly $C(r,s)$: the output coefficient of index $k$ is
$\sum_{i+j=k}x_i y_j$. The auxiliary rank-two identity contributes
the factor $B_X(2)$, leaving the coefficient leg unchanged.
For $F\in\mathcal V_{r-1}$, multiplication by $F$ respects evaluation
and its kernel. At input degree $s-1$ and output degree $r+s-2$,
its quotient and kernel blocks are multiplication
\[
 \mathcal V_s\longrightarrow\mathcal V_{r+s-1},\qquad
 \mathcal V_{s-2}\longrightarrow\mathcal V_{r+s-3}\quad(s\geq2).
\]
These are $C_{XY}(r,s+1)$ and $C_{XY}(r,s-1)$, sharing the
coefficient leg of $F$. Weight $-1$ on the input kernel and $+1$ on
the output dual kernel keeps both diagonal blocks at weight zero.
The block from input kernel to output quotient is zero, because
multiplication preserves the kernel. The opposite off-diagonal block
has weight one and disappears from the leading coefficient. Thus
$C_{XY}(r,s)\otimes B_X(2)$ degenerates onto these two matched
sectors, with a common first leg and disjoint second and third legs.
At $s=1$, the input kernel is empty, leaving only
$C_{XY}(r,2)$. Permuting legs gives the other orientations.

\paragraph{Three sectors.}
In $C(r,3s+r-1)$, the matched sectors have coordinates
($0\leq a<r$, $0\leq b<s$)
\[
 \begin{array}{c|ccc}
 &x&y&z\\\hline
 \text{left}&a&b&a+b\\
 \text{middle}&r-1-a&s+a+b&s+r-1+b\\
 \text{right}&a&2s+r-1+b&2s+r-1+a+b.
 \end{array}
\]
Give the middle $Y$-interval $[s,2s+r-2]$ weight $1$ and the
middle $Z$-interval $[s+r-1,2s+r-2]$ weight $-1$; all other weights
are zero. The convolution support prevents negative-weight terms.
The zero-weight tensor consists of the two outer $C_{XY}(r,s)$
sectors and the middle $C_{XZ}(r,s)$ sector, with the first
coordinate reversed in the middle sector.
That reversal identifies the value of the individual middle branch;
the actual three branches continue to share the original first leg.

\paragraph{Character profiles.}
The profiles are numerical functions of the two convolution dimensions.
A character may distinguish the first two legs. Pairing it with the
exchanged character makes one profile symmetric while retaining a second
profile for the different output orientation.
For a character $\lambda$ and its parameters from~\eqref{eq:dot},
define $G:\mathcal S\to\R_{\geq0}$ and $a_0,b_0$ by
\[
 G(T)=\sqrt{\lambda(T)\lambda(T^{X\leftrightarrow Y})},\qquad
 a_0=(p_X+p_Y)/2,\quad b_0=p_Z,
\]
Here $T^{X\leftrightarrow Y}$ exchanges the first two legs. For $r,s\geq1$,
define the symmetric and oriented convolution profiles
\[
 V(r,s)=G(C_{XY}(r,s)),\qquad U(r,s)=G(C_{XZ}(r,s)).
\]
The dot law, determinant endpoint, and output-dot restriction of
$C(2,2)$ give the bounds used at zero parameters:
\[
 V(r,1)=r^{a_0},\quad V(r,2)\leq2^{a_0}V(r,1),\quad
 2^{b_0}\leq V(2,2).
\]

For later leg permutations, write $\lambda^\pi(T)=\lambda(T^\pi)$,
where $XY,YZ$ exchange the named legs and
$T^{\mathrm{cyc}}(y,z,x)=T(x,y,z)$. The following identities follow by
relabeling the convolution indices. ``Paired'' means the geometric mean
of the two character values:
\[
\begin{array}{c|c|c|c|c}
\text{characters}&\text{arguments}&\text{paired }C_{XY}&
\text{paired }C_{XZ}&\text{shared exponent}\\\hline
(\lambda,\lambda^{XY})&(r,s)&V(r,s)&U(r,s)&a_0\\
(\lambda^{YZ},(\lambda^{XY})^{YZ})&(r,s)&U(r,s)&V(r,s)&a_0\\
(\lambda^{\mathrm{cyc}},(\lambda^{XY})^{\mathrm{cyc}})&
(s,r)&U(r,s)&U(r,s)&b_0
\end{array}
\]
For a matched-sector tensor $T=\sum_hT_h$, convert the sector
constructions into concavity by applying
Corollary~\ref{cor:entropy} to $\lambda$ and to the character obtained
by exchanging $X,Y$, using the same probability vector. Averaging
the logarithms gives, with averaged shared-leg exponent $\alpha>0$,
\[
 \log G(T)\geq\alpha H(q)+\sum_h q_h\log G(T_h).
\]
For positive sector values, put
$Z=\sum_hG(T_h)^{1/\alpha}$ and $q_h=G(T_h)^{1/\alpha}/Z$.
Then the right side is $\alpha\log Z$, yielding
\[
 G(T)^{1/\alpha}\geq\sum_h G(T_h)^{1/\alpha}
\]
when the averaged exponent of the shared leg is $\alpha>0$.
Apply this argument to each character pair in the table, using
degeneration monotonicity for the source. In the
determinant construction the source has the extra factor $B_X(2)$,
so its paired value is multiplied by $2^\alpha$. For a character pair
$(\phi,\psi)$, its paired convolution profile is
$J(r,s)=\sqrt{\phi(C_{XY}(r,s))\psi(C_{XY}(r,s))}$, with $\alpha$
the average of their $X$-dot exponents. Consequently,
\[
 J(r,s+1)^{1/\alpha}+J(r,s-1)^{1/\alpha}
 \leq2J(r,s)^{1/\alpha}\quad(s\geq2).
\]
In the third row, $J(r,s)=U(s,r)$, so this is first-coordinate
concavity for $U^{1/b_0}$. In the
three-sector construction there is no extra factor: the table gives
the two coupled tripling inequalities with exponent $a_0$, and the
self-tripling inequality with exponent $b_0$. Symmetry supplies the
other coordinate of $V$.

\section{The discrete parameter cap}

\begin{lemma}\label{lem:cap}
For every character, its dot parameters satisfy
\begin{equation}\label{eq:cap}
 b_0\leq4a_0(1-a_0),\qquad a_0=(p_X+p_Y)/2,\quad b_0=p_Z.
\end{equation}
\end{lemma}
\begin{proof}
Fix a character and its profiles just defined. If $a_0=0$, the
endpoint with $r=2$ gives $V(2,2)\leq1$, so $2^{b_0}\leq V(2,2)$
forces $b_0=0$. The case $b_0=0$ is immediate since $0\leq a_0\leq1$.
Hence assume $a_0,b_0>0$, and put
\[
 p=1/a_0\geq1,\qquad \theta=b_0/a_0>0,\qquad
 P=V^{1/a_0},\quad Q=U^{1/a_0}.
\]
The remaining goal is $\theta\leq4(p-1)/p$.
The determinant construction and its leg permutations make $P$
separately discretely concave, $Q$ concave in its second coordinate,
and $Q^{1/\theta}$ concave in its first. Here discrete concavity
means that consecutive increments are nonincreasing. Integer coordinate
restrictions and Lemma~\ref{lem:conv} give positivity, coordinate
monotonicity, symmetry of $P$, and
\[
 P(1,s)=s,\quad Q(n,1)=n^\theta,\qquad
 P(n,s),Q(n,s)\leq(n+s-1)^p\quad(n,s\geq1).
\]
The three-sector construction gives
\begin{align*}
 P(n,3s+n-1)&\geq2P(n,s)+Q(n,s),\\
 Q(n,3s+n-1)&\geq2Q(n,s)+P(n,s),\\
 Q(3n+s-1,s)^{1/\theta}&\geq3Q(n,s)^{1/\theta}.
\end{align*}
For $n\geq2$, the secant inequality for the concave row $P(n,\cdot)$ gives
\[
 P(n,4n-1)\leq P(n,n)+\frac{3n-1}{n-1}
                                (P(n,n)-P(n,1)).
\]
The first tripling inequality and positivity therefore give
$Q(n,n)\leq\frac{2n}{n-1}P(n,n)\leq4P(n,n)$.
A positive nondecreasing concave row satisfies $f(n)\leq nf(2)$: for $n\geq2$,
use $f(n)\leq f(1)+(n-1)(f(2)-f(1))$, and for $n=1$ use
monotonicity. Apply this in each coordinate of $P$ to obtain
\[
 n^\theta=Q(n,1)\leq Q(n,n)\leq4P(n,n)
                      \leq4P(2,2)n^2.
\]
Thus $\theta\leq2$. If $p\geq2$, this already gives
$\theta\leq4(p-1)/p$. It remains to handle $1\leq p<2$.

\paragraph{Common slopes.}
Row concavity gives limiting increments and the same linear growth
rates: the nonnegative increments decrease, and telescoping followed
by Ces\`aro averaging identifies their limit with $f(s)/s$.
For a fixed positive integer $n$, let $L_P,L_Q$ be these two row limits.
Dividing the first two tripling inequalities by $s$ gives
$3L_P\geq2L_P+L_Q$ and $3L_Q\geq2L_Q+L_P$, respectively,
so the two row growth rates agree. Write
\[
 L_n=\lim_{s\to\infty}\frac{P(n,s)}s
    =\lim_{s\to\infty}\frac{Q(n,s)}s.
\]
Every row increment is at least $L_n$, and monotonicity in $n$ makes
$L_n$ nondecreasing. Put $H(h)=P(n,h)+Q(n,h)$ for this paragraph.
Then $H(3h+n-1)\geq3H(h)$. Its iterates are
$h_j=3^j(h+(n-1)/2)-(n-1)/2$, so
$H(h)\leq H(h_j)/3^j$. Since $H(h_j)/h_j\to2L_n$, this yields
\[
 P(n,h)+Q(n,h)\leq(2h+n-1)L_n\quad(n,h\geq1).
\]
Positivity at $h=1$ implies $L_n>0$. With $D_n=P(n,n)$ and
$h=\max(n,m)$, use
$P(n,h)\geq D_n+(h-n)L_n$ and
$Q(n,h)\geq Q(n,m)+(h-m)L_n$ in the preceding bound to obtain
\[
 Q(n,m)\leq L_n(m+2n-1)-D_n.
\]
At $m=n$, $Q(n,n)\geq(n-1)L_n$ gives $D_n\leq2nL_n$.
Also $D_n\geq n+(n-1)L_n$, since $P(n,1)=n$. For $n\geq2$,
these imply
\[
 \tfrac12\leq\frac{nL_n}{D_n}\leq2\quad(n\geq2).
\]

\paragraph{Eliminate the varying slope.}
Fix positive integers $u,v$, let $\tau=u/v$, and, for integers $t\geq2$,
set
\[
 D(t)=P(vt,vt),\quad W(t)=Q(vt,ut),\quad
 \ell_t=\frac{vtL_{vt}}{D(t)},\qquad z=\frac{u-1}{2v+u}\geq0.
\]
For positive integers $n,m$, the two-coordinate gains follow in a
definite order. First, symmetry
and the nondecreasing slopes give
$P(n+v,n+v)\geq D_n+vL_n+vL_{n+v}\geq D_n+2vL_n$.
For $Q$, move in the second coordinate first. The affine bound gives
\[
 Q(n,m+u)\geq Q(n,m)+uL_n
 \geq Q(n,m)\left(1+
          \frac{uL_n}{L_n(m+2n)-D_n}\right).
\]
The denominator is positive because it strictly exceeds $Q(n,m)$.
At fixed second coordinate $s$, put $f(j)=Q(j,s)^{1/\theta}$.
Concavity and self-tripling give, for $1\leq h\leq2n+s-1$,
\[
 f(n+h)-f(n)\geq\frac{h}{2n+s-1}(f(3n+s-1)-f(n))
             \geq\frac{2h}{2n+s-1}f(n).
\]
Use $h=v$, $n=vt$, $m=ut$, and then $s=m+u$. The resulting steps are
\begin{align*}
 D(t+1)&\geq D(t)(1+2\ell_t/t),\\
 W(t+1)&\geq W(t)(1+F_\tau(\ell_t)/t)
                       (1+\gamma/(t+z))^\theta,\\
 F_\tau(x)&=\frac{\tau x}{(\tau+2)x-1},\qquad
 \gamma=\frac2{2+\tau}.
\end{align*}
The function $F_\tau$ is used on $x\geq1/2$, where its denominator is positive.
The varying $\ell_t\in[1/2,2]$ is the obstacle to summing the
second step directly. To eliminate it, put $y=p/2$. The identity
\[
 F_\tau(x)-F_\tau(y)+
 \frac{\tau(x-y)}{((\tau+2)y-1)^2}
 =\frac{\tau(\tau+2)(x-y)^2}
 {((\tau+2)x-1)((\tau+2)y-1)^2}\geq0
\]
for $x,y\geq1/2$ proves the tangent bound
\[
 F_\tau(x)\geq F_\tau(p/2)-c(x-p/2),\qquad
 c=\frac{\tau}{((\tau+2)p/2-1)^2}\geq0\quad(x\geq1/2).
\]
The tangent contributes $-c\ell_t/t$, which is cancelled by
$(c/2)$ times the diagonal logarithmic increment. Accordingly, the
potential $\Phi(t)=\log W(t)+(c/2)\log D(t)$ satisfies
\begin{align*}
 \Phi(t+1)-\Phi(t)&\geq
 \frac{\theta\gamma+F_\tau(p/2)+cp/2}{t}-O(t^{-2}),\\
 \Phi(t)&\leq p(1+c/2)\log t+O(1).
\end{align*}
The first estimate uses $\log(1+x)\geq x-x^2/2$; the bounded slopes
make the error uniform in $t$ for this fixed ray. The second uses the
rank bounds. Summing telescopically, the $t^{-2}$ errors remain bounded
and the $cp/2$ terms cancel, leaving
\[
 \frac{2\theta}{2+\tau}+\frac{\tau}{\tau+2-2/p}\leq p.
\]
This holds for every positive rational $\tau$, hence for every
$\tau>0$ by continuity. Taking $\tau=4/p-2>0$ gives
$\theta\leq4(p-1)/p$. Since $b_0=a_0\theta$ and $p=1/a_0$,
this is~\eqref{eq:cap}.
\end{proof}

\section{The finite scheme and its arithmetic cost}

\begin{proof}[Proof of Theorem~\ref{thm:ring}]
Fix $k\geq0$ and $\varepsilon>0$, and put $\eta=\varepsilon/2$.
Choose an integer $q\geq2$ with $2<q^\eta$, and let
$j=\lceil q^k\rceil$. For every character, $j\leq2q^k$ and
$0\leq b_0\leq1$, so~\eqref{eq:mm} and Lemma~\ref{lem:cap} give
\[
 \lambda(\langle q,j,q\rangle)
 \leq2q^{2a_0+kb_0}
 \leq2q^{\max_{0\leq x\leq1}(2x+4kx(1-x))}=2q^{f(k)}.
\]
The quadratic maximum occurs at $x=1$ for $k\leq1/2$, and at
$x=1/2+1/(4k)$ otherwise. Hence~\eqref{eq:spectrum} gives
\[
 \ark(\langle q,j,q\rangle)\leq2q^{f(k)}<q^{f(k)+\eta}.
\]
There is therefore a fixed integer $m\geq1$ with an exact integer
rank decomposition of length $\rho<q^{m(f(k)+\eta)}$. Reindexing
this tensor power gives $\langle a,b,a\rangle$, where
$a=q^m$, $b=j^m$. Its three flattenings over $\Q$ supply the last
condition in the finite base certificate
\begin{equation}\label{eq:scheme}
 a\geq2,\qquad b\geq a^k,\qquad
 \log_a\rho<f(k)+\eta,\qquad \rho\geq\max\{ab,a^2\}.
\end{equation}
Only this fixed scheme is needed from now on.

\paragraph{Correctness over every ring.}
Let $R$ be any associative ring, and let
$A\in R^{a\times b}$ and $B\in R^{b\times a}$.
Writing the chosen integer decomposition coefficients as
$\alpha,\beta,\gamma$ gives, for $1\leq i,h\leq a$,
\[
 (AB)_{ih}=\sum_{\ell=1}^{\rho}\gamma_{ih,\ell}
 \left(\sum_{i',j'}\alpha_{\ell,i'j'}A_{i'j'}\right)
 \left(\sum_{j'',h'}\beta_{\ell,j''h'}B_{j''h'}\right),
\]
where the entry sums run over the indicated matrix coordinates.
Expansion leaves each $A$-entry to the left of its $B$-entry, and
integer coefficient equality proves the identity. Integer multiples
use only addition and negation. The same expansion applies to
rectangular block multiplication, a $\Z$-bilinear map, so it also
justifies the recursive scheme.

\paragraph{Arithmetic cost and padding.}
An integer coefficient $c$ uses at most $|c|+1$ gates, and adding
its term to an accumulator uses one further gate. Let $K_0$ be the
sum of $|c|+2$ over all entries of the three fixed coefficient arrays,
and set $d=\max\{ab,a^2\}\leq\rho$. The preparation and reconstruction
at depth $t+1$ act on blocks with $(ab)^t,(ab)^t,(a^2)^t$ entries.
Each of the $\rho$ recursive products is computed once and reused in
reconstruction. Thus the actual gate count $T(t)$ satisfies
\[
 T(0)=1,\qquad T(t+1)\leq\rho T(t)+K_0d^t.
\]
Induction using $d\leq\rho$ gives
\[
 \operatorname{cost}(t)=T(t)\leq(1+K_0)(t+1)\rho^t.
\]
For an integer $n\geq1$, choose $t=\lfloor\log_a n\rfloor+1$.
Then
\[
 n<a^t\leq an,\qquad \lceil n^k\rceil\leq b^t.
\]
Zero input padding and output wires give the requested shape.
The inner padding need not have a constant ratio to $n^k$; its
operations are already included in $\operatorname{cost}(t)$.
Since $\log_a\rho\geq0$ and $t+1=O(1+\log n)$,
\[
 \operatorname{cost}(t)=O\bigl((1+\log n)n^{\log_a\rho}\bigr)
                     =O_{k,\varepsilon}(n^{f(k)+\varepsilon}).
\]
Here $1+\log n=O_\eta(n^\eta)$ and~\eqref{eq:scheme} reserves the
other $\eta$ of slack. A fixed constant covers every $n\geq1$.
The base dimensions, coefficients, and operation counts depend only
on $k,\varepsilon$, so this constant is independent of $R$.
\end{proof}

\section{Supplementary exponent bounds}

For nonnegative shape exponents $u,v,w$, let
$\omega_{\mathrm{univ}}(u,v,w)$ be the infimum of arithmetic exponents
of universal integer bilinear circuits for shapes
$\lceil n^u\rceil\times\lceil n^v\rceil$ and
$\lceil n^v\rceil\times\lceil n^w\rceil$. Let $\omega_R(u,v,w)$ be
its counterpart when algorithms may depend on the fixed ring $R$.
The main theorem gives
\[
 \omega_R(1,k,1)\leq\omega_{\mathrm{univ}}(1,k,1)\leq f(k).
\]
For the universal lower bound, evaluate the integer coefficient
identities over $\Q$. Output flattening gives at least $n^2$
products, hence
$\omega_{\mathrm{univ}}(1,k,1)=2$ for $0\leq k\leq1/2$.
In particular the universal dual exponent is at least $1/2$.
A fixed degenerate ring can have weaker arithmetic lower bounds;
counting explicit dense output writes also gives exponent two for
that output model.

\begin{corollary}[Three arbitrary dimensions]\label{cor:three}
Let $a\geq b\geq c\geq0$ be the decreasing rearrangement of three
shape exponents. Uniformly over all associative rings,
\[
 \omega_R(a,b,c)\leq\omega_{\mathrm{univ}}(a,b,c)
 \leq
 \begin{cases}
 a+b,&b\geq2c,\\[2pt]
 a+b+\dfrac{(2c-b)^2}{4(a+c-b)},&b<2c.
 \end{cases}
\]
The bound holds for every permutation of the shape exponents,
and the universal exponent equals $a+b$ when $b\geq2c$.
\end{corollary}
\begin{proof}
Permute~\eqref{eq:cap} to obtain the three symmetric parameter
constraints. These permutations act on integer tensor identities;
each resulting scheme still places the first input on the left and
the second on the right.
For any character, let $(p_1,p_2,p_3)$ be the permutation of its dot
parameters corresponding to the sorted shape weights, and put
$S=p_1+p_2+p_3$. If $S\leq2$, then $p_i\in[0,1]$ bounds the
objective $ap_1+bp_2+cp_3$ by $a+b$.
If $S>2$, the three cap constraints imply $S\leq9/4$ and, with
$d=\sqrt{9-4S}\in[0,1]$,
\[
 \frac{3-2d-d^2}{4}\leq p_i\leq\frac{3+2d-d^2}{4}.
\]
Using these endpoints and $p_1+p_2+p_3=S$ bounds the objective by
\[
 \frac{3(a+b+c)}4+\frac{a-c}2d-\frac{a+c-b}4d^2.
\]
Its maximum on $[0,1]$ is the displayed bound; in the second branch,
$b<2c$ ensures $a+c-b>0$.

For rational shape exponents, choose a common integer base $A\geq2$
so $A^a,A^b,A^c$ are integers. The character bound and spectral
duality produce a finite power scheme with rank ratio below the
claimed bound plus any positive slack. At common depth
$t=\lfloor\log_{A^m}n\rfloor+1$, where $m\geq1$ is that fixed power,
the same level estimate gives $O((t+1)\rho^t)$ operations.
For real exponents, approximate the sorted triple from above by
rational triples and use continuity of the bound and padding.
When all exponents are zero, a single scalar product suffices.
Finally, the three flattenings over $\Q$ give the universal lower
bound $\max\{a+b,a+c,b+c\}=a+b$.
\end{proof}

\paragraph{Verification.}
The accompanying \texttt{verify\_integral\_constructions.py} checks
finite instances of the separator, determinant bases and blocks,
and the oriented tripling identity using exact integer arithmetic.
These checks supplement the general proofs. The \texttt{formalization}
directory proves Theorem~\ref{thm:ring}, including coefficient arithmetic,
with strict addition, negation, and ordered multiplication syntax.
The witnesses are classical circuit existence proofs; no executable
implementation or evaluator runtime bound is extracted. The
three-dimension corollary and lower bounds are supplementary mathematical
statements outside the public Lean theorem.

\begin{thebibliography}{9}
\bibitem{BCS} Peter B\"urgisser, Michael Clausen, and M. Amin Shokrollahi.
\emph{Algebraic Complexity Theory}. Springer, 1997.
Chapters 14--15, pp.~351--423.
\url{https://doi.org/10.1007/978-3-662-03338-8}.
\bibitem{LN} Rudolf Lidl and Harald Niederreiter.
\emph{Finite Fields}. Cambridge University Press.
Chapter 2, ``Structure of Finite Fields.''
\url{https://doi.org/10.1017/CBO9780511525926}.
\bibitem{U} Przemys\l aw Uzna\'nski.
\emph{Rectangular matrix multiplication from shared-leg entropy}.
arXiv:2610.10494v1, October 7, 2026.
\url{https://arxiv.org/html/2610.10494v1}.
\bibitem{O} OpenAI.
\emph{An upper bound of $9/4$ for the matrix multiplication exponent}.
October 2, 2026.
\url{https://github.com/openai/math/tree/main/preprints/Matrix-Multiplication-Nine-Fourths-October-2-2026}.
\bibitem{Zui} Jeroen Zuiddam.
\emph{Algebraic complexity, asymptotic spectra and entanglement polytopes}.
Updated dissertation, April 6, 2020. Sections 2.3, 2.7--2.9.
\url{https://staff.fnwi.uva.nl/j.zuiddam/diss/diss_v2.pdf}.
\end{thebibliography}
\end{document}
